%==============================================================================
%  Chapter: The PESAD System
%  Section: Domain Analysis
%==============================================================================
\chapter{The PESAD System}

The following sections present the analysis of the domain, the phases of
identification, conceptualization and knowledge representation and, finally, the
implementation of the functional components of the realized expert system.

\section{Domain Analysis}

This section discusses the domain under study for the construction of the expert
system. A brief introduction to the concept of ``anxiety'' is followed by its
distinction into ``state anxiety'' and ``trait anxiety''. The symptoms of
anxiety are then introduced, together with their classification according to the
functional system they affect. Finally, the classification adopted in the
present work is described: that of the fifth edition, text revision, of the
\emph{Diagnostic and Statistical Manual of Mental Disorders}
(DSM-5-TR, 2022)~\cite{dsm}, which distributes the conditions of interest across
three distinct chapters, each of which is analysed in turn.

\subsection{Anxiety}

From an etymological point of view, anxiety derives from the Latin term
\emph{anxia} (to tighten). Anxiety is a universal emotion which, in itself, it
would not be inappropriate to experience, since it represents a necessary part
of the response to stress. It constitutes, in fact, a defence mechanism aimed at
anticipating the perception of danger even before it has clearly manifested
itself. Anxiety typically sets in motion the physiological mechanisms that push,
on the one hand, towards exploration in order to identify the danger and confront
it in the most appropriate manner and, on the other hand, towards avoidance and
flight.

Anxiety is generally regarded by mental-health professionals as ``the mother of
all emotions'' and, precisely for this reason, it is useful to differentiate an
anxiety understood as an existential condition (``trait'' anxiety) from a
pathological anxiety (``state'' anxiety).

The distinction between the concepts of ``trait'' anxiety and ``state'' anxiety
was introduced by Cattell and Scheier (1961) and further elaborated by
Spielberger and colleagues in the development of their self-assessment scale, the
State--Trait Anxiety Inventory~\cite{stai}. The two paragraphs below report some
salient points of the two concepts of anxiety, as elaborated by Spielberger and
colleagues.

\subsubsection{Trait Anxiety}

Trait anxiety can be considered a relatively stable characteristic of the
personality, a behavioural attitude, which reflects the manner in which the
subject tends to perceive environmental stimuli and situations as dangerous or
threatening.

Subjects with high trait anxiety show a more marked reactivity to a greater
number of stimuli and are characterized, according to Cattell and Scheier, by
high arousal, ego weakness, and a tendency towards sensitivity and guilt.

These characteristics seem to indicate a sort of predisposition to anxiety, in
the sense that subjects with high trait anxiety have a greater probability, with
respect to others, of presenting state anxiety in circumstances with low
anxiety-provoking potential and/or, for equal stimuli, of experiencing higher
levels of state anxiety. For subjects with high trait anxiety, the anxious one is
the habitual mode of response to stimuli and to environmental situations, whereas
for the others it is an exceptional mode.

\subsubsection{State Anxiety}

State anxiety manifests itself as an interruption of the emotional continuum,
that is, it provokes a rupture in the emotional balance of the person; it is
expressed through a subjective sensation of tension, worry, restlessness,
nervousness and reactivity. It is associated with an activation of the autonomic
nervous system, which provokes a series of physiological activations.

High levels of state anxiety are particularly unpleasant, disturbing and even
painful, to the point of inducing the person to put into action behavioural
adaptation mechanisms aimed at putting an end to these sensations. However, these
mechanisms may fail to achieve their purpose, leaving room for other behaviours,
this time of a maladaptive type, which lead to the opposite effect, namely a
further increase of anxiety, thus starting a vicious circle of a pathological
nature.

State anxiety, therefore, beyond certain levels of severity and duration, becomes
a pathology that may take the form of an \emph{anxiety disorder} or of one of the
related conditions described in the remainder of this chapter.

In this work, the anxiety under study is pathological state anxiety, that is, the
state anxiety that becomes disproportionate to the events that arouse it, that
persists even when the events have been overcome, or that manifests itself in the
absence of objective events.

\subsection{Symptoms of Anxiety}

The symptoms that denote an anxiety disorder may generally involve one or more of
the four functional systems that are normally coordinated by our organism to
produce adaptive responses to dangerous situations: the cognitive system, the
affective system, the behavioural system and the physiological system.

Generally, in the presence of a threat, the organism reacts by activating the
different systems: the cognitive apparatus establishes the existence of a danger
and provides selective evaluations of the threat and of the resources available
to face it; the affective component acts as an accelerator of the reaction by
increasing the sense of urgency; the behavioural component activates or inhibits
action; and the physiological component acts as a somatic mobilizer.

\paragraph{Cognitive symptoms.}
The cognitive symptoms include three different types of symptom: sensory
symptoms, difficulties of thought, and conceptual symptoms.

\begin{itemize}[leftmargin=2em]
  \item \textbf{Sensory symptoms} are significant not so much because of their
        severity or their interference with sensory and perceptual function, but
        because, being little known and not immediately controllable, they tend
        to induce more sensitive people to think they are about to ``lose their
        mind''. Such symptoms include: a confused or clouded mind, visualization
        of objects in a blurred/distant way, visualization of the environment in
        a different/unreal way, sense of unreality, and hypervigilance.
  \item \textbf{Difficulties of thought} mainly concern difficulty in
        remembering important things, confusion, the inability to control one's
        thinking, difficulty of concentration, distractibility, blocking,
        difficulty in reasoning, and loss of objectivity and perspective. They
        may be produced by a variety of factors, among them the cognitive
        inhibitions that interfere with memory and the continuous coping with
        danger.
  \item \textbf{Conceptual symptoms} reflect the preoccupation with the sense of
        vulnerability and danger. Examples are the fear of losing control, the
        fear of not being able to cope with situations, the fear of physical
        injury or of death, the fear of mental disorders, the fear of negative
        evaluations, threatening visual images, and repetitive frightening
        ideation.
\end{itemize}

\paragraph{Affective symptoms.}
The affective symptoms are often the most conspicuous in anxiety disorders. The
qualitative experiences of anxiety may differ from one situation to another, and
even from time to time within the same situation, so the anxious person may often
feel irritable, impatient, ill at ease, nervous, tense, touchy, fearful,
frightened, terrified, alarmed, appalled, excited, agitated, and so on.

\paragraph{Behavioural symptoms.}
The behavioural symptoms of anxiety generally reflect either the hyperactivity of
the behavioural system or its inhibition. There may indeed be inhibition,
immobility of muscle tone, flight, avoidance, difficult speech, defective
coordination, agitation, collapse and hyperventilation.

\paragraph{Physiological symptoms.}
The physiological symptoms of anxiety reflect the readiness of the organism, in
its globality, to react in order to protect itself by activating the branch of
the autonomic nervous system. The physiological symptoms in the different anxiety
disorders may involve the functional systems summarized in
Table~\ref{tab:physiological}.

\begin{longtable}{@{}p{0.30\textwidth}p{0.62\textwidth}@{}}
  \caption{Physiological symptoms of anxiety, grouped by the bodily system they
  affect.}\label{tab:physiological}\\
  \toprule
  \textbf{System / activity} & \textbf{Symptoms} \\
  \midrule
  \endfirsthead
  \multicolumn{2}{@{}l}{\small\itshape Table~\ref{tab:physiological} (continued)}\\
  \toprule
  \textbf{System / activity} & \textbf{Symptoms} \\
  \midrule
  \endhead
  \midrule
  \multicolumn{2}{r}{\small\itshape continued on next page}\\
  \endfoot
  \bottomrule
  \endlastfoot
  Cardiovascular activity &
    palpitations, increased heart rate, increased blood pressure, weakness,
    fainting, decreased blood pressure, decreased heart rate \\
  Respiratory activity &
    breathing difficulties, pressure in the chest, lump in the throat, sensation
    of choking, laboured, rapid or shallow breathing, etc. \\
  Neuromuscular activity &
    increased reflexes, startle reaction, twitching eyelid, insomnia, spasm,
    tremor, stiffness, agitation, strained facial expression, restless pacing,
    swaying, generalized weakness, wobbly legs, clumsy movements \\
  Gastrointestinal activity &
    abdominal pain, loss of appetite, repulsion towards food, nausea, heartburn,
    abdominal discomfort, vomiting \\
  Urinary-tract activity &
    urge to urinate, frequency of urination \\
  Skin activity &
    flushing of the face or pallid face, localized or diffuse sweating, hot or
    cold flushes, itching \\
\end{longtable}

\subsection{The DSM-5-TR Classification}

The reference nosology adopted in this work is the fifth edition, text revision,
of the \emph{Diagnostic and Statistical Manual of Mental Disorders}
(DSM-5-TR), published by the American Psychiatric Association in
2022~\cite{dsm}. With respect to the earlier DSM-IV-TR, on which the first version
of PESAD was based, the DSM-5-TR introduced a substantial reorganization of the
conditions historically grouped together as ``anxiety disorders''. The single
DSM-IV-TR class is now split across \textbf{three} separate chapters:

\begin{itemize}[leftmargin=2em]
  \item \textbf{Anxiety Disorders};
  \item \textbf{Obsessive-Compulsive and Related Disorders}; and
  \item \textbf{Trauma- and Stressor-Related Disorders}.
\end{itemize}

The key reorganizations introduced by this split, and relevant to the present
system, may be summarized as follows.

\begin{itemize}[leftmargin=2em]
  \item Obsessive-Compulsive Disorder is no longer classified among the anxiety
        disorders; it heads the new \emph{Obsessive-Compulsive and Related
        Disorders} chapter, together with body dysmorphic disorder, hoarding
        disorder, trichotillomania and excoriation disorder.
  \item Posttraumatic Stress Disorder and Acute Stress Disorder are likewise no
        longer anxiety disorders; they belong to the new \emph{Trauma- and
        Stressor-Related Disorders} chapter, which also requires an identifiable
        traumatic or stressful event as part of its diagnostic criteria.
  \item Panic disorder and agoraphobia, formerly coded jointly (as ``panic
        disorder with/without agoraphobia'' and ``agoraphobia without history of
        panic disorder''), are now \textbf{two separate diagnoses}, each of which
        can be coded independently and both of which may be assigned together.
        The \emph{panic attack} itself is no longer a codable disorder but a
        \textbf{cross-cutting specifier} that can be appended to any diagnosis
        (e.g.\ ``PTSD with panic attacks'').
  \item Social phobia was renamed \textbf{social anxiety disorder}.
  \item \textbf{Hoarding disorder} and \textbf{excoriation (skin-picking)
        disorder} are entirely new diagnoses, not present in DSM-IV-TR.
  \item \textbf{Prolonged grief disorder} is a new diagnosis added specifically
        in the DSM-5-TR text revision of 2022; it was not present in DSM-IV-TR,
        nor in the original DSM-5 of 2013.
\end{itemize}

The following three subsections describe, chapter by chapter, the disorders
actually modelled by PESAD. For each disorder a single paragraph summarizes its
core DSM-5-TR criteria, namely the key symptoms, the polythetic thresholds (of
the form ``at least $M$ of $N$'') and the required duration. The residual
\emph{other specified} and \emph{unspecified} categories of each chapter are not
modelled as positive diagnostic goals and are therefore omitted.

\subsection{Anxiety Disorders}

\subsubsection{Separation Anxiety Disorder}

The disorder is defined by developmentally inappropriate and excessive fear or
anxiety concerning separation from those to whom the individual is attached, as
evidenced by \textbf{at least 3 of 8} symptoms: recurrent excessive distress when
anticipating or experiencing separation; persistent worry about losing
attachment figures or about harm befalling them; persistent worry that an
untoward event will cause separation; reluctance or refusal to go out for fear of
separation; fear of being alone; reluctance or refusal to sleep away from home or
without a major attachment figure nearby; repeated nightmares about separation;
and repeated complaints of physical symptoms when separation occurs or is
anticipated. The fear, anxiety or avoidance must be persistent, lasting
\textbf{at least 4 weeks in children and adolescents} (younger than 18 years) and
\textbf{typically 6 months or more in adults}, and must cause clinically
significant distress or impairment.

\subsubsection{Selective Mutism}

Selective mutism is characterized by a consistent failure to speak in specific
social situations in which there is an expectation to speak (e.g.\ at school),
despite speaking in other situations. The disturbance must interfere with
educational or occupational achievement or with social communication, last
\textbf{at least 1 month} (not limited to the first month of school), and not be
attributable to a lack of knowledge of, or comfort with, the spoken language
required. It is not better explained by a communication disorder and does not
occur exclusively during an autism spectrum or psychotic disorder.

\subsubsection{Specific Phobia}

Specific phobia involves marked fear or anxiety about a specific object or
situation, which almost always provokes immediate fear and is actively avoided or
endured with intense anxiety. The fear is out of proportion to the actual danger
and to the sociocultural context, is persistent (\textbf{typically 6 months or
more}, at all ages) and causes clinically significant distress or impairment. The
disorder is coded according to the phobic stimulus, and PESAD models the
\textbf{five} DSM-5-TR stimulus types listed in Table~\ref{tab:phobia-types}.

\begin{longtable}{@{}p{0.34\textwidth}p{0.58\textwidth}@{}}
  \caption{The five stimulus types of specific phobia.}%
  \label{tab:phobia-types}\\
  \toprule
  \textbf{Type} & \textbf{Typical stimuli} \\
  \midrule
  \endfirsthead
  \toprule
  \textbf{Type} & \textbf{Typical stimuli} \\
  \midrule
  \endhead
  \bottomrule
  \endlastfoot
  Animal & spiders, insects, dogs \\
  Natural environment & heights, storms, water \\
  Blood--injection--injury & needles, invasive medical procedures, injury \\
  Situational & aeroplanes, lifts, enclosed places \\
  Other & choking or vomiting situations; in children, loud sounds or costumed
          characters \\
\end{longtable}

\subsubsection{Social Anxiety Disorder}

Social anxiety disorder is defined by marked fear or anxiety about one or more
social situations involving possible scrutiny by others (social interactions,
being observed, or performing), driven by the fear of acting in a way, or showing
anxiety symptoms, that will be negatively evaluated. The social situations almost
always provoke fear, are avoided or endured with intense anxiety, and the fear is
out of proportion to the actual threat. The disturbance is persistent
(\textbf{typically 6 months or more}) and causes clinically significant distress
or impairment. A \emph{performance only} specifier applies when the fear is
restricted to speaking or performing in public.

\subsubsection{Panic Disorder}

Panic disorder requires recurrent \textbf{unexpected} panic attacks. A panic
attack is an abrupt surge of intense fear or discomfort that peaks within
minutes, during which \textbf{at least 4 of the 13} symptoms of
Table~\ref{tab:panic-symptoms} occur. At least one attack must have been followed
by \textbf{1 month or more} of persistent concern about further attacks or their
consequences, or of significant maladaptive change in behaviour related to the
attacks. The disturbance must not be attributable to a substance or another
medical condition, nor better explained by another mental disorder.

\begin{longtable}{@{}p{0.05\textwidth}p{0.88\textwidth}@{}}
  \caption{The thirteen symptoms of a panic attack (eleven physical and two
  cognitive); at least four are required.}\label{tab:panic-symptoms}\\
  \toprule
  \# & \textbf{Symptom} \\
  \midrule
  \endfirsthead
  \toprule
  \# & \textbf{Symptom} \\
  \midrule
  \endhead
  \bottomrule
  \endlastfoot
  1 & Palpitations, pounding heart, or accelerated heart rate \\
  2 & Sweating \\
  3 & Trembling or shaking \\
  4 & Shortness of breath or smothering \\
  5 & Feelings of choking \\
  6 & Chest pain or discomfort \\
  7 & Nausea or abdominal distress \\
  8 & Feeling dizzy, unsteady, light-headed, or faint \\
  9 & Chills or heat sensations \\
  10 & Paraesthesias (numbness or tingling) \\
  11 & Derealization or depersonalization \\
  12 & Fear of losing control or ``going crazy'' \\
  13 & Fear of dying \\
\end{longtable}

\subsubsection{Agoraphobia}

Agoraphobia is defined by marked fear or anxiety about \textbf{at least 2 of the
5} situations: using public transportation; being in open spaces; being in
enclosed places; standing in line or being in a crowd; and being outside the home
alone. The situations are feared or avoided because of thoughts that escape might
be difficult, or help unavailable, should panic-like or other incapacitating or
embarrassing symptoms develop. They almost always provoke fear and are actively
avoided, require a companion, or are endured with intense anxiety. The fear is out
of proportion to the actual danger, is persistent (\textbf{typically 6 months or
more}) and causes clinically significant distress or impairment. Agoraphobia is
diagnosed irrespective of panic disorder; if both sets of criteria are met, both
diagnoses are assigned.

\subsubsection{Generalized Anxiety Disorder}

Generalized anxiety disorder is characterized by excessive anxiety and worry,
occurring \textbf{more days than not for at least 6 months}, about a number of
events or activities, which the individual finds difficult to control. The
anxiety and worry must be associated with \textbf{at least 3 of the 6} symptoms
(only \textbf{1} required in children): restlessness or feeling keyed up;
being easily fatigued; difficulty concentrating or mind going blank;
irritability; muscle tension; and sleep disturbance. The disturbance causes
clinically significant distress or impairment and is not attributable to a
substance or another medical condition.

\subsubsection{Substance/Medication-Induced Anxiety Disorder}

In this disorder, panic attacks or anxiety predominate in the clinical picture,
and there is evidence that the symptoms developed during or soon after substance
intoxication or withdrawal, or after exposure to or withdrawal from a medication
capable of producing such symptoms. The disturbance must not be better explained
by an independent anxiety disorder; persistence of symptoms for a substantial
period (\textbf{about 1 month}) after the cessation of acute withdrawal or severe
intoxication suggests an independent disorder. The symptoms must not occur
exclusively during a delirium and must cause clinically significant distress or
impairment. The condition is coded with onset during intoxication, during
withdrawal, or after medication use.

\subsubsection{Anxiety Disorder Due to Another Medical Condition}

Here panic attacks or anxiety predominate in the clinical picture, and there is
evidence from the history, physical examination or laboratory findings that the
disturbance is the direct pathophysiological consequence of another medical
condition (e.g.\ hyperthyroidism, hyperparathyroidism, phaeochromocytoma,
vestibular dysfunction, seizure disorders, or cardiopulmonary conditions). The
disturbance must not be better explained by another mental disorder, must not
occur exclusively during a delirium, and must cause clinically significant
distress or impairment.

\subsection{Obsessive-Compulsive and Related Disorders}

\subsubsection{Obsessive-Compulsive Disorder}

Obsessive-compulsive disorder is defined by the presence of obsessions,
compulsions, or both. \emph{Obsessions} are recurrent and persistent thoughts,
urges or images that are experienced as intrusive and unwanted and that cause
marked anxiety or distress, which the individual attempts to ignore, suppress or
neutralize with another thought or action. \emph{Compulsions} are repetitive
behaviours or mental acts that the individual feels driven to perform in response
to an obsession or according to rigidly applied rules, aimed at preventing or
reducing distress but not connected in a realistic way with what they are
designed to neutralize, or clearly excessive. The obsessions or compulsions must
be time-consuming (\textbf{taking more than 1 hour per day}) or cause clinically
significant distress or impairment. An \textbf{insight} specifier with three
levels (good or fair, poor, absent or delusional) and a \emph{tic-related}
specifier are available.

\subsubsection{Body Dysmorphic Disorder}

Body dysmorphic disorder consists of a preoccupation with one or more perceived
defects or flaws in physical appearance that are not observable or appear slight
to others. At some point the individual must have performed repetitive behaviours
(e.g.\ mirror checking, excessive grooming, skin picking, reassurance seeking) or
mental acts (e.g.\ comparing one's appearance with that of others) in response to
the appearance concerns. The preoccupation must cause clinically significant
distress or impairment and must not be better explained by concerns with body fat
or weight in an individual whose symptoms meet criteria for an eating disorder. A
\emph{with muscle dysmorphia} specifier and the same three-level insight specifier
apply.

\subsubsection{Hoarding Disorder}

Hoarding disorder is characterized by a persistent difficulty discarding or
parting with possessions, regardless of their actual value, owing to a perceived
need to save the items and to the distress associated with discarding them. This
difficulty results in the accumulation of possessions that congest and clutter
active living areas and substantially compromise their intended use. The hoarding
must cause clinically significant distress or impairment (including the
maintenance of a safe environment) and must not be attributable to another
medical condition or better explained by another mental disorder. A \emph{with
excessive acquisition} specifier and the three-level insight specifier apply.

\subsubsection{Trichotillomania (Hair-Pulling Disorder)}

Trichotillomania is defined by recurrent pulling out of one's hair resulting in
hair loss, together with repeated attempts to decrease or stop the pulling. The
hair pulling must cause clinically significant distress or impairment, must not
be attributable to another medical condition (e.g.\ a dermatological condition),
and must not be better explained by another mental disorder (e.g.\ attempts to
improve a perceived appearance defect in body dysmorphic disorder).

\subsubsection{Excoriation (Skin-Picking) Disorder}

Excoriation disorder consists of recurrent skin picking resulting in skin
lesions, together with repeated attempts to decrease or stop the picking. The
skin picking must cause clinically significant distress or impairment, must not
be attributable to the physiological effects of a substance (e.g.\ cocaine) or
another medical condition (e.g.\ scabies), and must not be better explained by
another mental disorder.

\subsubsection{Substance/Medication-Induced Obsessive-Compulsive and Related
Disorder}

In this disorder, obsessions, compulsions, skin picking, hair pulling, other
body-focused repetitive behaviours or other symptoms characteristic of the
chapter predominate in the clinical picture, and there is evidence that the
symptoms developed during or soon after substance intoxication or withdrawal, or
after exposure to or withdrawal from a medication capable of producing such
symptoms. The disturbance must not be better explained by an independent
obsessive-compulsive and related disorder, must not occur exclusively during a
delirium, and must cause clinically significant distress or impairment. It is
coded with onset during intoxication, during withdrawal, or after medication use.

\subsubsection{Obsessive-Compulsive and Related Disorder Due to Another Medical
Condition}

Here obsessions, compulsions, appearance preoccupations, hoarding, hair pulling,
skin picking or other characteristic symptoms predominate in the clinical
picture, and there is evidence that the disturbance is the direct
pathophysiological consequence of another medical condition. The disturbance must
not be better explained by another mental disorder, must not occur exclusively
during a delirium, and must cause clinically significant distress or impairment.
Specifiers identify the predominant symptom presentation (obsessive-compulsive,
appearance preoccupations, hoarding, hair-pulling, or skin-picking).

\subsection{Trauma- and Stressor-Related Disorders}

\subsubsection{Reactive Attachment Disorder}

Reactive attachment disorder is defined by a consistent pattern of inhibited,
emotionally withdrawn behaviour towards adult caregivers, manifested by
\textbf{both} of: the child rarely or minimally seeking comfort when distressed,
and rarely or minimally responding to comfort when distressed. It is accompanied
by a persistent social and emotional disturbance shown by \textbf{at least 2 of
3} symptoms (minimal social and emotional responsiveness; limited positive
affect; episodes of unexplained irritability, sadness or fearfulness). The child
must have experienced a pattern of extremes of insufficient care
(\textbf{at least 1 of 3}: social neglect, repeated changes of caregivers, or
rearing in settings limiting selective attachments), presumed responsible for the
disturbance. The criteria for autism spectrum disorder must not be met, the
disturbance must be evident \textbf{before age 5}, and the child must have a
developmental age of \textbf{at least 9 months}.

\subsubsection{Disinhibited Social Engagement Disorder}

This disorder is defined by a pattern of behaviour in which the child actively
approaches and interacts with unfamiliar adults, shown by \textbf{at least 2 of
4} symptoms: reduced or absent reticence in approaching unfamiliar adults; overly
familiar verbal or physical behaviour; diminished or absent checking back with a
caregiver after venturing away; and willingness to go off with an unfamiliar
adult with minimal hesitation. The behaviours must extend beyond mere
impulsivity, must follow a pattern of extremes of insufficient care
(\textbf{at least 1 of 3}, as for reactive attachment disorder) presumed
responsible for them, and the child must have a developmental age of
\textbf{at least 9 months}.

\subsubsection{Posttraumatic Stress Disorder}

Posttraumatic stress disorder requires exposure to actual or threatened death,
serious injury or sexual violence through one or more of four routes (directly
experiencing the event; witnessing it in person; learning that it happened to a
close family member or friend; or experiencing repeated or extreme exposure to
aversive details). The symptoms are organized into the four clusters of
Table~\ref{tab:ptsd-clusters}; the disturbance must last \textbf{more than 1
month} and cause clinically significant distress or impairment. A \emph{with
dissociative symptoms} specifier (depersonalization or derealization) and a
\emph{with delayed expression} specifier (full criteria not met until at least 6
months after the event) are available.

\begin{longtable}{@{}p{0.20\textwidth}p{0.10\textwidth}p{0.60\textwidth}@{}}
  \caption{The four symptom clusters of PTSD (individuals older than 6 years) and
  their polythetic thresholds.}\label{tab:ptsd-clusters}\\
  \toprule
  \textbf{Cluster} & \textbf{Threshold} & \textbf{Content} \\
  \midrule
  \endfirsthead
  \toprule
  \textbf{Cluster} & \textbf{Threshold} & \textbf{Content} \\
  \midrule
  \endhead
  \bottomrule
  \endlastfoot
  B -- Intrusion & at least 1 of 5 &
    intrusive memories; distressing dreams; dissociative reactions
    (flashbacks); psychological distress at, and physiological reactivity to,
    cues resembling the event \\
  C -- Avoidance & at least 1 of 2 &
    avoidance of distressing memories, thoughts or feelings; avoidance of
    external reminders of the event \\
  D -- Negative cognitions and mood & at least 2 of 7 &
    amnesia for the event; exaggerated negative beliefs; distorted blame;
    persistent negative emotional state; diminished interest; detachment from
    others; inability to feel positive emotions \\
  E -- Arousal and reactivity & at least 2 of 6 &
    irritable or aggressive behaviour; reckless or self-destructive behaviour;
    hypervigilance; exaggerated startle; concentration problems; sleep
    disturbance \\
\end{longtable}

For \textbf{children 6 years and younger} (the preschool subtype) the thresholds
are developmentally lowered: exposure occurs through 3 routes rather than 4, the
intrusion cluster keeps the threshold of at least 1 of 5, and avoidance and
negative cognitions are merged into a single cluster requiring \textbf{at least 1
of 6} symptoms, while the arousal cluster requires \textbf{at least 2 of 5}. The
duration (more than 1 month) and the available specifiers are unchanged.

\subsubsection{Acute Stress Disorder}

Acute stress disorder shares the same exposure criterion as PTSD (one or more of
four routes). It requires the presence of \textbf{at least 9 of 14} symptoms
drawn from five categories (intrusion, negative mood, dissociation, avoidance and
arousal), beginning or worsening after the event. The duration of the disturbance
is \textbf{3 days to 1 month after the trauma}; symptoms resolving in fewer than 3
days do not qualify, and symptoms persisting beyond 1 month lead instead to a
diagnosis of PTSD. The disturbance must cause clinically significant distress or
impairment, must not be attributable to a substance or another medical condition,
and must not be better explained by a brief psychotic disorder.

\subsubsection{Adjustment Disorder}

Adjustment disorder is the development of emotional or behavioural symptoms in
response to an identifiable stressor, occurring \textbf{within 3 months} of the
onset of the stressor. The symptoms must be clinically significant, shown by
marked distress out of proportion to the stressor and/or significant impairment.
The disturbance must not meet the criteria for another mental disorder, must not
be merely an exacerbation of a pre-existing disorder, and must not represent
normal bereavement or be better explained by prolonged grief disorder. Once the
stressor or its consequences have terminated, the symptoms do not persist for
more than a further \textbf{6 months}. The subtype is specified as: with depressed
mood; with anxiety; with mixed anxiety and depressed mood; with disturbance of
conduct; with mixed disturbance of emotions and conduct; or unspecified. Course
specifiers distinguish \emph{acute} (less than 6 months) from \emph{persistent}
(6 months or longer).

\subsubsection{Prolonged Grief Disorder}

Prolonged grief disorder, new in DSM-5-TR, requires the death of a close person
\textbf{at least 12 months ago} (\textbf{at least 6 months} in children and
adolescents). Since the death there must be a persistent grief response shown by
\textbf{one or both} of intense yearning or longing for the deceased and
preoccupation with thoughts or memories of the deceased, present nearly every day
for at least the last month. In addition, \textbf{at least 3 of 8} symptoms must
be present nearly every day for at least the last month: identity disruption; a
marked sense of disbelief about the death; avoidance of reminders that the person
is dead; intense emotional pain related to the death; difficulty reintegrating
into relationships and activities; emotional numbness; a feeling that life is
meaningless; and intense loneliness. The disturbance must cause clinically
significant distress or impairment, must clearly exceed expected social, cultural
or religious norms, and must not be better explained by another mental disorder
nor attributable to a substance or another medical condition.
